% =====================================================================
%  Part II -- The physical substrate.
%  Source: tier0.md Part 0 preamble (123-137), 0.1 (138-306),
%          0.2 (308-439), 0.3 (441-548), 0.4 (550-641).
% =====================================================================
\part{The physical substrate}
\label{part:substrate}

The physics that must be in place before ``a nuclide has $(Z,N)$'' and ``the
network is $\dd\mathbf{Y}/\dd t = \nu\mathbf{R}$'' mean anything.

\begin{center}
\small
\begin{tabularx}{\textwidth}{@{}l R{0.680} R{1.320}@{}}
\toprule
\textbf{\S} & \textbf{Area} &
\textbf{Why it is prerequisite, not optional} \\
\midrule
\ref{sec:thermostate} & Plasma thermodynamics \& EOS &
  Fixes the regime box, determines \emph{which} screening prescription is
  valid, sets whether electron capture is even energetically allowed \\
\ref{sec:statmech} & Statistical mechanics &
  NSE, partition functions, and detailed balance are one framework; without it
  S3 and S9 are memorized formulae \\
\ref{sec:nuclearstructure} & Nuclear structure \& data &
  Mass excess vs binding energy is a real sign trap; and most rates here are
  \emph{theoretical}, not measured \\
\ref{sec:neutrinos} & Neutrino physics &
  Explains the acceleration of late burning, the net one-way \Ye{} ratchet, and
  an open project item ($\langle E_\nu\rangle$) \\
\bottomrule
\end{tabularx}
\end{center}

Numbers below are computed at box-centre conditions and are worth re-deriving
yourself --- several justify configuration choices that otherwise look
arbitrary.

% ---------------------------------------------------------------------
\chapter{The thermodynamic state of silicon-burning matter}
\label{sec:thermostate}

\section{The four pressure components}
\label{sec:pressurecomponents}

\begin{equation*}
P = P_{\mathrm{ion}} + P_{e} + P_{\mathrm{rad}} + P_{\mathrm{Coulomb}}
\end{equation*}

\begin{itemize}
\item \textbf{Ions} --- ideal gas, $P_{\mathrm{ion}} = n_i kT$,
      $n_i = \rho N_A/\bar{A}$
\item \textbf{Electrons} --- degenerate, relativistic; the dominant term here
\item \textbf{Radiation} --- $P_{\mathrm{rad}} = aT^4/3$,
      $a = 7.566\times10^{-15}$~erg\,cm$^{-3}$\,K$^{-4}$
      \citep{codata2022}
\item \textbf{Coulomb} --- negative correction from ion--ion and ion--electron
      correlation; small for pressure, but \emph{this is the term that becomes
      screening} (\S\ref{sec:gamma})
\end{itemize}

\section{Worked example --- box centre
  (\texorpdfstring{$\rho = 10^{8}$\,\gcc, $T_9 = 4$, $Y_e = 0.5$,
   $\bar{A} \approx 28$}{rho = 1e8, T9 = 4, Ye = 0.5, Abar = 28})}
\label{sec:boxcentre}

\begin{align*}
n_e &= \rho N_A Y_e = 10^{8}\times 6.022\times10^{23}\times 0.5
       = 3.01\times10^{31}\ \mathrm{cm^{-3}}\\
n_i &= \rho N_A/\bar{A} = 10^{8}\times 6.022\times10^{23}/28
       = 2.15\times10^{30}\ \mathrm{cm^{-3}}\\
kT  &= 8.617\times10^{-2}\ \mathrm{MeV\,GK^{-1}}\times 4\ \mathrm{GK}
       = 0.345\ \mathrm{MeV}
\end{align*}

Fermi momentum from $n_e$ (degenerate, $T \to 0$ limit):

\begin{align*}
p_F c &= \hbar c\,(3\pi^2 n_e)^{1/3}, \qquad
         \hbar c = 1.9733\times10^{-11}\ \mathrm{MeV\,cm}\\
&= 1.9733\times10^{-11}\times(29.608\times 3.01\times10^{31})^{1/3}\\
&= 1.9733\times10^{-11}\times 9.62\times10^{10}\\
&= 1.90\ \mathrm{MeV}
\end{align*}

Relativity parameter $x = p_F c / m_e c^2 = 1.90 / 0.511 = \mathbf{3.7}$ ---
relativistic. Total Fermi energy including rest mass:

\begin{equation*}
E_F = \sqrt{p_F^2c^2 + m_e^2c^4} = \sqrt{3.60+0.26} = 1.97\ \mathrm{MeV}
\end{equation*}

Pressures (ultra-relativistic degenerate approximation
$P_e \approx \tfrac{1}{4} n_e p_F c$):

\begin{center}
\begin{tabular}{@{}l r r@{}}
\toprule
\textbf{Component} & \textbf{Value (erg/cm$^3$)} & \textbf{Share} \\
\midrule
$P_e$              & $2.3\times10^{25}$ & \textbf{${\sim}93\%$} \\
$P_{\mathrm{ion}}$ & $1.2\times10^{24}$ & ${\sim}5\%$ \\
$P_{\mathrm{rad}}$ & $6.5\times10^{23}$ & ${\sim}3\%$ \\
\bottomrule
\end{tabular}
\end{center}

\textbf{Electrons carry the pressure.} This is the quantitative content of
\S\ref{sec:explodability} --- the core is an electron-degeneracy-supported
object, so a composition variable that changes $n_e$ changes the structure
directly.

\section{Degeneracy across the box}
\label{sec:degeneracy}

Degeneracy is measured by $E_{F,\mathrm{kin}}/kT$, with
$E_{F,\mathrm{kin}} = E_F - m_ec^2$:

\begin{center}
\small
\begin{tabularx}{\textwidth}{@{}l l l l Y@{}}
\toprule
\textbf{$\rho$ (\gcc)} & \textbf{$p_F c$ (MeV)} &
\textbf{$E_{F,\mathrm{kin}}$ (MeV)} &
\textbf{$E_{F,\mathrm{kin}}/kT$ at $\Tn = 4$} & \textbf{Regime} \\
\midrule
$10^{7}$ & 0.88 & 0.51 & 1.5  & mildly degenerate \\
$10^{8}$ & 1.90 & 1.46 & 4.2  & degenerate \\
$10^{9}$ & 4.09 & 3.61 & 10.5 & strongly degenerate \\
\bottomrule
\end{tabularx}
\end{center}

The box spans exactly the transition from ``electrons barely notice each
other'' to ``electrons dominate everything'' --- not a coincidence, but the
physically interesting range.

\section{The Coulomb coupling parameter
  \texorpdfstring{$\Gamma$}{Gamma} --- and why \code{chugunov\_2007}}
\label{sec:gamma}

Ion-sphere radius and coupling parameter:

\begin{equation*}
a_i = \left(\frac{3}{4\pi n_i}\right)^{1/3}, \qquad
\Gamma = \frac{Z^2 e^2}{a_i kT}, \qquad
e^2 = 1.44\times10^{-13}\ \mathrm{MeV\,cm}
\end{equation*}

At $\rho = 10^{8}$, $\Tn = 4$:

\begin{equation*}
a_i = \left(\frac{3}{4\pi\times 2.15\times10^{30}}\right)^{1/3}
    = 4.81\times10^{-11}\ \mathrm{cm}
\end{equation*}

\begin{align*}
\text{Silicon }(Z=14):\quad \Gamma &=
  \frac{196\times 1.44\times10^{-13}}{4.81\times10^{-11}\times 0.345} = 1.7\\
\text{Iron }(Z=26):\quad \Gamma &=
  \frac{676\times 1.44\times10^{-13}}{6.05\times10^{-11}\times 0.345} = 4.7
\end{align*}

Now scale it across the box. $\Gamma \propto n_i^{1/3}/T \propto \rho^{1/3}/T$,
so \textbf{both} axes of the box move it, and the temperature axis moves it by
more:

\begin{center}
\begin{tabular}{@{}l c c@{}}
\toprule
 & \textbf{Si ($Z=14$, $\bar{A}=28$)} & \textbf{Fe ($Z=26$, $\bar{A}=56$)} \\
\midrule
$\rho = 10^{7}$, $\Tn = 7.9$ & \textbf{0.40} & 1.10 \\
$\rho = 10^{7}$, $\Tn = 4$   & 0.79 & 2.16 \\
$\rho = 10^{8}$, $\Tn = 4$   & 1.70 & 4.66 \\
$\rho = 10^{9}$, $\Tn = 4$   & 3.67 & 10.05 \\
$\rho = 10^{9}$, $\Tn = 1.6$ & 9.18 & \textbf{25.1} \\
\bottomrule
\end{tabular}
\end{center}

So across the \textbf{full} box ($\rho = 10^{7}$--$10^{9}$ \emph{and}
$T = 1.6$--7.9~GK, Si $\to$ Fe) \textbf{$\Gamma$ ranges roughly 0.4 to 25}. The
often-quoted ``0.8 to 10'' is the $\rho$-only span at fixed $\Tn = 4$ ---
correct as far as it goes, but it silently freezes the axis that matters most.
Quote the wider range.

\begin{aside}
\textbf{This is the derivation that justifies the screening config.} Salpeter
weak screening \citep{salpeter1954} assumes $\Gamma \ll 1$. Strong-screening
asymptotics assume $\Gamma \gg 1$. Silicon burning sits in \textbf{neither} ---
and the full-box range makes the point more forcefully than the fixed-$T$ one
does, because $\Gamma \sim 0.4$ to ${\sim}25$ is two orders of magnitude
straddling unity, with no sub-region where either asymptotic is safe. You need
a prescription that \emph{interpolates across both limits}. That is exactly
what \code{chugunov\_2007} is --- the fit of \citet{chugunov2007} reproduces the
Debye--Hückel $\Gamma \ll 1$ asymptote and is calibrated through
$\Gamma \sim$ hundreds, spanning this box --- and it is why the label
configuration pins \code{screening\_mode = 'chugunov'} in
\code{scripts/bbq\_campaign/inlist.template}.
\end{aside}

Carry this into \textbf{S4}: when you read \code{fluxes/screening.py}, you are
reading an intermediate-coupling interpolation, and you now know why nothing
simpler would do.

\section{The electron-capture threshold --- a box edge derived from scratch}
\label{sec:ecthreshold}

Electron capture on free protons, $p + e^- \to n + \nu$, requires the total
electron energy to exceed the neutron--proton mass difference
\citep{codata2022}:

\begin{equation*}
E_e \geq (m_n-m_p)c^2 = 1.2933\ \mathrm{MeV}
\qquad (\text{kinetic threshold } 0.782\ \mathrm{MeV})
\end{equation*}

In a degenerate gas the \emph{available} electron energy is $E_F$. Setting
$E_F = 1.2933$~MeV:

\begin{align*}
p_F c &= \sqrt{1.2933^2 - 0.511^2} = 1.188\ \mathrm{MeV}\\
n_e &= \frac{(p_F c)^3}{3\pi^2(\hbar c)^3}
     = \frac{1.6765}{29.608\times 7.684\times10^{-33}}
     = 7.37\times10^{30}\ \mathrm{cm^{-3}}\\
\rho &= \frac{n_e}{N_A Y_e}
      = \frac{7.37\times10^{30}}{6.022\times10^{23}\times 0.5}
      = 2.4\times10^{7}\ \gcc
\end{align*}

Sanity check: at $\Ye = 1$ this gives $1.2\times10^{7}$~\gcc{} --- the classic
neutronization threshold for hydrogen \citep{shapiro1983}; captures on H switch
on once $\rho \gtrsim 10^{7}$~\gcc{} \citep{bildsten1998}. \yes{}

\begin{result}
\textbf{Result:} free-proton electron capture switches on at
$\rho \approx 2.4\times10^{7}$~\gcc{} for $\Ye = 0.5$ --- \emph{inside} the
regime box, near its lower edge.
\end{result}

Below that density captures are thermally suppressed (only the distribution's
tail has enough energy); above it they are degeneracy-driven and vigorous.
\textbf{The box's lower density edge of $10^{7}$~\gcc{} sits just below the
threshold on purpose} --- it brackets the turn-on of the process that controls
\Ye{}.

Captures on nuclei have their own thresholds set by their $Q$-values and in
practice carry most of the capture flux --- heavy nuclei outnumber free protons
by orders of magnitude at these conditions \citep{langanke2003}; the
free-proton number is the clean reference point.

\begin{selfcheck}
\item Recompute $p_F c$, $E_F$, and $\Gamma$ at $\rho = 10^{9}$, $\Tn = 3$, and
      confirm degeneracy and coupling both increase. Which pressure component
      is largest?
\item Show $P_e \propto (\rho \Ye)^{4/3}$ in the ultra-relativistic limit, and
      connect it to the $M_{\mathrm{ch}} \propto Y_e^2$ derivation in
      \S\ref{sec:mchderivation}.
\item At what temperature would radiation pressure overtake ion pressure at
      $\rho = 10^{8}$? Is that inside the box?
\item Why does $\Gamma$ \emph{increase} as burning proceeds at fixed
      $(\rho, T)$? (Hint: what happens to $\bar{Z}$?) What does that imply
      about screening corrections drifting along a trajectory?
\end{selfcheck}

% ---------------------------------------------------------------------
\chapter{Statistical mechanics you actually need}
\label{sec:statmech}

Three results, one framework. Everything in S3 (detailed balance) and S9 (NSE
and QSE) is a corollary of this chapter.

\section{Chemical potential and the equilibrium condition}
\label{sec:chempotential}

For a system at fixed $(T, V)$, the Helmholtz free energy $F$ is minimized. The
chemical potential of species $i$ is

\begin{equation*}
\mu_i = \left(\frac{\partial F}{\partial N_i}\right)_{T,V,N_{j\neq i}}
\end{equation*}

For a reaction with stoichiometric coefficients $\nu_i$ (products $+$,
reactants $-$), advancing by $\dd\xi$ changes $N_i$ by $\nu_i \dd\xi$, so
$\dd F = \sum_i \mu_i \nu_i \dd\xi$. Equilibrium ($\dd F = 0$ for arbitrary
$\dd\xi$) gives the \textbf{chemical equilibrium condition}:

\begin{equation}
\sum_i \nu_i \mu_i = 0
\label{eq:chem-equilibrium}
\end{equation}

This single equation generates everything:

\begin{itemize}
\item \textbf{NSE.} Every strong/EM reaction is equilibrated. Apply
      \eqref{eq:chem-equilibrium} to the build-up of nuclide $(Z,N)$ from free
      nucleons, $A(Z,N) \rightleftharpoons Z\,p + N\,n$:
      %
      \begin{equation}
      \mu(Z,N) = Z\mu_p + N\mu_n
      \label{eq:nse}
      \end{equation}
      %
      Every nuclide's chemical potential is determined by just \textbf{two}
      numbers --- the standard NSE condition
      \citep{hix1996,hix1999b,pitrou2018}. That is why the NSE solver in
      \code{qse/solver.py} has exactly two unknowns ($u_p$, $u_n$).

\item \textbf{QSE.} Only a \emph{subset} of reactions is equilibrated --- those
      internal to a cluster. Cluster members satisfy \eqref{eq:nse} up to a
      shared offset:
      %
      \begin{equation}
      \mu_i = Z_i\mu_p + N_i\mu_n + u_G \quad \text{for } i\in G
      \label{eq:qse-offset}
      \end{equation}
      %
      One extra unknown per cluster --- the focal-nucleus form of
      \citet{hix1996}, whose $r_{\mathrm{QSE}}$ is constant within a group and
      offset between groups, and whose QSE distribution ``blends simply into
      the NSE distribution'' as the offset vanishes. That is the third unknown
      in \code{solve\_qse}, and $u_G = 0$ recovers NSE exactly --- which is why
      the QSE$\to$NSE degeneracy is a test case in
      \code{tests/test\_qse\_solver.py}.

\item \textbf{Detailed balance.} Apply \eqref{eq:chem-equilibrium} to a single
      forward/reverse pair; the ratio of rate coefficients is fixed by the
      $Q$-value and partition functions. This is S3.
\end{itemize}

\section{Nuclear partition functions}
\label{sec:partitionfunctions}

A nucleus has excited states. The internal partition function is

\begin{equation}
G(T) = \frac{\sum_{\text{states}} (2J_s+1)\,e^{-E_s/kT}}{2J_0+1}
\label{eq:partition-fn}
\end{equation}

normalized so $G \to 1$ as $T \to 0$ (only the ground state populated).

\textbf{Why it matters here, quantitatively.} First excited states in
Si--Fe-peak nuclei span a wide range (\citeauthor{ensdf}): tens of keV for odd-$A$ and
odd--odd species (\nuc{Fe}{57} at 14.4~keV, \nuc{Mn}{54} at 54.9~keV),
${\sim}0.8$--1.2~MeV for even--even Fe-group nuclei (\nuc{Fe}{56} at
0.847~MeV), up to 1.8--2.7~MeV for the Si-shell even--even and near-magic cases
(\nuc{Si}{28} at 1.78~MeV, \nuc{Ni}{56} at 2.70~MeV). At $\Tn = 4$,
$kT = 0.345$~MeV, so a ${\sim}0.85$~MeV state carries
$e^{-2.5} \approx 0.08$ per magnetic substate while states below ${\sim}kT$ are
populated at order unity. $G$ reaches ${\sim}1.5$--3 for typical odd-$A$
Fe-group nuclei --- and up to ${\sim}4$--7 for the extreme low-lying cases
(\nuc{Fe}{57} $\approx 6.6$ at $\Tn = 4$) --- while the doubly-magic and
near-magic Fe-peak nuclei (\nuc{Ni}{56} above all, $G = 1.002$ at $\Tn = 4$)
stay at $G \approx 1$ \citep{rauscher2000,rauscher2003}. So the pf correction
is strongly species-dependent, and largest exactly where the level density is,
not where the abundance is. (``Deformation'' is the wrong word for this regime;
near $N = Z$ and $Z = 28$ these nuclei are close to spherical. What raises $G$
here is level density at $E^* \sim kT$.)

\begin{aside}
A factor of 2 in $G$ is a factor of 2 in a detailed-balance reverse rate. That
is the entire reason \code{DerivedRate(use\_pf=True)} is mandatory above
$\Tn \approx 3$, and why raw pf-free fits manufacture spurious $\kappa$ floors.
\end{aside}

\begin{repopointer}
\textbf{In the code:} partition functions come from Rauscher tables
\citep[][--- the reverse-rate ratio is linear in $G$, so a factor of 2 in $G$
is exactly a factor of 2 in the reverse rate]{rauscher2000,rauscher2003},
rebuilt as splines with constant extrapolation, in both
\code{fluxes/compile.py} and \code{qse/coeffs.py} --- and the two paths are
cross-checked against each other, because a pf mismatch between the rate side
and the equilibrium side would silently corrupt every $\kappa$ comparison.
\end{repopointer}

\section{Detailed balance from time reversal}
\label{sec:detailedbalance}

The microscopic statement: for a transition $a \to b$, the squared matrix
element is invariant under time reversal --- $T$-invariance implies $S$-matrix
symmetry $S_{ab} = S_{ba}$ and hence detailed balance
$|S_{ab}|^2 = |S_{ba}|^2$ \citep{mitchell2010} ---

\begin{equation*}
\bigl|\langle b|T|a\rangle\bigr|^2 = \bigl|\langle a|T|b\rangle\bigr|^2
\end{equation*}

The \emph{rates} differ only through phase-space density of final states and
statistical weights. In equilibrium, forward and reverse rates per unit volume
must be equal --- that is what ``detailed'' balance means: not just net zero
overall, but \textbf{pairwise} zero, reaction by reaction.

This pairwise statement is exactly what $\kappa$ measures:

\begin{equation}
\kappa_r = \frac{|f^+ - f^-|}{f^+ + f^-} \;\longrightarrow\; 0
\quad \text{at true equilibrium}
\label{eq:kappa-detailed-balance}
\end{equation}

So $\kappa$ is not an arbitrary diagnostic --- it is the numerical test of
detailed balance, reaction by reaction. A nonzero $\kappa$ floor at NSE means
the code's forward and reverse rates are \emph{mutually inconsistent}, which is
a bug, not physics. That is the logic of the whole $\kappa$-floor screen.
(\S\ref{sec:cancellation} gives $\kappa$ its second, independent meaning as a
condition number --- reaching the identical formula
\eqref{eq:kappa-condition} from numerical analysis rather than thermodynamics.
The two readings coincide, which is why $\kappa$ is load-bearing everywhere in
this project.)

\begin{selfcheck}
\item Derive \eqref{eq:nse} from \eqref{eq:chem-equilibrium} for the specific
      case $\nuc{Si}{28} \rightleftharpoons 14p + 14n$.
\item Explain why NSE needs exactly two constraints ($\sum X = 1$, \Ye{}) given
      \eqref{eq:nse} has two free parameters. What breaks if you impose a
      third?
\item Estimate $G(T)$ at $\Tn = 5$ for a nucleus with first excited state at
      0.85~MeV and $J = 2$ (ground $J = 0$). By what factor would omitting it
      change a detailed-balance reverse rate?
\item Argue why detailed balance is \emph{pairwise} and not merely a statement
      about net flow. Then explain how a code could satisfy overall equilibrium
      while having every individual $\kappa_r \neq 0$.
\end{selfcheck}

% ---------------------------------------------------------------------
\chapter{Nuclear structure minimum}
\label{sec:nuclearstructure}

\section{The mass-excess convention --- a real sign trap}
\label{sec:massexcess}

Nuclear data tables give the \textbf{mass excess}, not the binding energy:

\begin{equation}
\Delta(A,Z) \equiv \bigl[M_{\text{atomic}}(A,Z) - A\,u\bigr]\,c^2
\label{eq:mass-excess}
\end{equation}

Tabulated masses are \textbf{atomic} (electrons included). Binding energy is
recovered by (mass excesses throughout from AME2020 ---
\citealt{ame2020}):

\begin{equation}
\begin{aligned}
B(A,Z) &= Z\,\Delta(\nuc{H}{1}) + N\,\Delta(n) - \Delta(A,Z),\\
\Delta(\nuc{H}{1}) &= 7.2890\ \mathrm{MeV}, \qquad
\Delta(n) = 8.0713\ \mathrm{MeV}
\end{aligned}
\label{eq:binding}
\end{equation}

\textbf{Worked example --- \nuc{Fe}{56}} ($\Delta = -60.607$~MeV):

\begin{align*}
B &= 26\times 7.2890 + 30\times 8.0713 + 60.607\\
&= 189.51 + 242.14 + 60.61 = 492.26\ \mathrm{MeV}\\
B/A &= 8.790\ \mathrm{MeV} \qquad \yes\ \text{textbook value}
\end{align*}

\begin{center}
\begin{tabular}{@{}l r l@{}}
\toprule
\textbf{Nuclide} & \textbf{$\Delta$ (MeV)} & \textbf{$B/A$ (MeV)} \\
\midrule
\nuc{Ni}{56} ($Z=N=28$) & $-53.908$ & 8.643 \\
\nuc{Fe}{56}            & $-60.607$ & \textbf{8.790} \\
\nuc{Fe}{58}            & $-62.155$ & 8.792 \\
\nuc{Ni}{62}            & $-66.746$ & \textbf{8.795} $\leftarrow$ most bound
                                                                nuclide \\
\bottomrule
\end{tabular}
\end{center}

(\nuc{Ni}{62} as the most bound nuclide --- not \nuc{Fe}{56}, a common
misconception --- is \citealt{fewell1995}: $B/A = 8794.60 \pm 0.03$~keV.)

\textbf{$Q$-values follow directly:}

\begin{equation}
Q = \sum \Delta_{\text{reactants}} - \sum \Delta_{\text{products}}
\label{eq:qvalue}
\end{equation}

\textbf{Worked example --- $\nuc{Si}{28}(\alpha,\gamma)\nuc{S}{32}$:}

\begin{align*}
Q &= \bigl[\Delta(\nuc{Si}{28}) + \Delta(\nuc{He}{4})\bigr]
     - \Delta(\nuc{S}{32})\\
&= [-21.493 + 2.425] - (-26.016)\\
&= -19.068 + 26.016 = 6.948\ \mathrm{MeV} \qquad \yes
\end{align*}

This matters concretely: \code{configs/reactions\_mesa\{80,151\}\_mesa.yaml}
carries a \code{q\_mev} field per reaction, and the energy identity
$\enuc = \sum_j Q_j \Phi_j$ in S10 consumes it. You should be able to reproduce
any entry by hand.

\section{Magic numbers and why \texorpdfstring{\nuc{Ni}{56}}{56Ni} is the
  Si-burning endpoint}
\label{sec:magicnumbers}

Shell closures at $N$ or $Z = 2, 8, 20, 28, 50, 82, 126$ produce extra binding
\citep{otsuka2020}. \textbf{$Z = N = 28$ makes \nuc{Ni}{56} doubly magic},
which is why:

\begin{itemize}
\item it is the most bound nuclide accessible at $\Ye = 0.5$,
\item it is the dominant Si-burning product \citep{hix1996},
\item and its decay chain \nuc{Ni}{56} ($T_{1/2} = 6.08$~d) $\to$ \nuc{Co}{56}
      ($T_{1/2} = 77.2$~d) $\to$ \nuc{Fe}{56} powers supernova light curves
      (\citealt{diehl2015}; half-lives from \citeauthor{ensdf}), which is how we
      \emph{measure} how much of it was made (\S\ref{sec:obsgrounding}).
\end{itemize}

Note the mechanism connecting to \S\ref{sec:explodability}: the binding-energy
maximum (\nuc{Ni}{62}) and the actual product (\nuc{Ni}{56}) differ because the
\textbf{\Ye{} constraint} binds. Weak reactions are the only way to relax it,
and they are slow. The composition is set by a \emph{constrained} optimum, and
the constraint is the project's central variable.

\section{Where the rates actually come from --- and how well they are known}
\label{sec:ratesources}

The part most easily assumed away. In a network reaching to Ni and beyond:

\begin{center}
\small
\begin{tabularx}{\textwidth}{@{}R{0.75} R{1.05} R{1.20}@{}}
\toprule
\textbf{Source} & \textbf{Coverage} & \textbf{Typical uncertainty} \\
\midrule
Direct experiment &
  light nuclei, a minority of channels &
  a few \% (best-measured channels) to a few tens of \%
  \citep{iliadis2010} \\
Hauser--Feshbach statistical model &
  \textbf{the large majority} of $(n,\gamma)$, $(p,\gamma)$, $(\alpha,\gamma)$
  on mid/heavy nuclei &
  factor ${\sim}1.5$--2, worse off stability \citep{rauscher2000} \\
Shell-model weak rates (LMP) &
  Fe-peak electron capture / $\beta$-decay &
  no blanket figure; GT-strength dependent --- the LMP revision moved key FFN
  EC rates by factors ${\sim}10$--350 \citep{langanke2000} \\
\bottomrule
\end{tabularx}
\end{center}

REACLIB is a \emph{fit library} \citep{cyburt2010}: seven coefficients per set,
several sets summed per rate, fitted to whichever of the above is best
available. Two consequences:

\begin{enumerate}
\item \textbf{The seven-coefficient form is a fitting basis with physical
      asymptotics, not an exact theory} (S2 derives why those particular
      terms).
\item \textbf{There is an irreducible physics-uncertainty floor.} An emulator
      reproducing MESA to $10^{-6}$ is reproducing \emph{MESA}, not nature.
      This is why the project frames its gates as fidelity-to-labels --- and
      why the label-pathology finding (S13) matters so much: it is the case
      where MESA and nature part company by 7--13 dex.
\end{enumerate}

\begin{selfcheck}
\item Reproduce $B/A$ for \nuc{Ni}{56} from \eqref{eq:binding} and confirm
      8.643~MeV.
\item Compute $Q$ for $\nuc{S}{32}(\alpha,\gamma)\nuc{Ar}{36}$ and for the
      reverse. What is the sign relation?
\item Pick five reactions from \code{configs/reactions\_mesa80\_mesa.yaml} and
      verify \code{q\_mev} by hand from a mass table.
\item Given that most $(n,\gamma)$ rates carry factor-${\sim}2$ theoretical
      uncertainty, what is the strongest \emph{defensible} claim an emulator
      can make about physical accuracy? Write it as one sentence.
\end{selfcheck}

% ---------------------------------------------------------------------
\chapter{Neutrinos}
\label{sec:neutrinos}

\section{Two distinct neutrino channels}
\label{sec:nuchannels}

Do not conflate these; the dataset carries them differently.

\textbf{(a) Thermal neutrino losses} --- pair annihilation
($e^+e^- \to \nu\bar\nu$), photoneutrino, plasmon decay, bremsstrahlung
\citep{itoh1996,kato2020}. Depend on the thermodynamic state ($T$, $\rho \Ye$)
for the $e^\pm$ processes, plus mean $\bar{Z}/\bar{A}$ for bremsstrahlung ---
not on the detailed isotopic abundances. Pair production dominates at
silicon-burning temperatures \citep{kato2020} and approaches its
relativistic-limit $\propto T^9$ scaling there (steeper, ${\sim}T^{11}$ or
more, at lower $T$).

\textbf{(b) Weak-reaction neutrinos} --- emitted by every electron capture and
$\beta$ decay. Depend on composition, carry away lepton number as well as
energy, and are inseparable from the \Ye{} evolution.

In the data: \code{eps\_nu} / \code{eps\_neu} columns. Note the units trap from
\S\ref{sec:labelnoise} --- a \textbf{rate} [erg/g/s] in the training CSVs while
\code{eps\_nuc} beside it is \textbf{integrated} [erg/g].

\section{Why neutrinos free-stream --- and why that accelerates everything}
\label{sec:freestreaming}

Neutrino--nucleon cross sections scale as \citep[][where
$\sigma_0 = (4/\pi)G_F^2(m_ec^2)^2/(\hbar c)^4 = 1.761\times10^{-44}$~cm$^2$]
{janka2017}

\begin{equation*}
\sigma \approx \sigma_0\left(\frac{E_\nu}{m_e c^2}\right)^{\!2}, \qquad
\sigma_0 \approx 1.7\times10^{-44}\ \mathrm{cm^2}
\end{equation*}

At $E_\nu \approx 3$~MeV: $\sigma \approx 5.9\times10^{-43}$~cm$^2$. At
$\rho = 10^{8}$~\gcc, $n_{\mathrm{baryon}} = 6.0\times10^{31}$~cm$^{-3}$:

\begin{equation*}
\lambda = \frac{1}{n\sigma}
        = \frac{1}{6.0\times10^{31}\times 5.9\times10^{-43}}
        \approx 2.8\times10^{10}\ \mathrm{cm}
        \approx 2.8\times10^{5}\ \mathrm{km}
\end{equation*}

The iron core radius is ${\sim}10^{3}$~km, so
$\boldsymbol{\lambda/R \approx 300}$: neutrinos escape without scattering.
(Coherent scattering on nuclei enhances $\sigma$ by roughly
$N^2/16 \approx A^2/60$ --- the vector amplitudes on the $N$ neutrons add in
phase \citep{janka2017} --- cutting $\lambda$, but trapping only sets in around
$\rho \approx 10^{11}$--$10^{12}$~\gcc{} \citep{janka2017}, well above this
box.)

\textbf{Two consequences, both structural:}

\begin{enumerate}
\item \textbf{Energy leaves instantly.} No diffusion time. This is the
      mechanism behind the burning-stage acceleration table in
      \S\ref{sec:lastday} --- the core must burn ever faster to replace energy
      leaking at $c$.
\item \textbf{Lepton number leaves with it.} Each capture emits a $\nu$ that
      never comes back. This is why the constraint matrix needs a
      \emph{separate neutrino ledger}: lepton number is conserved by the
      reaction but then physically removed from the zone.

      \begin{warn}
      \textbf{Be careful with the next step, which is easy to overstate.}
      Free-streaming does not itself \emph{drive} \Ye{} down --- what drives it
      down is electron capture outpacing $\beta^-$ decay, a statement about net
      rates and about how degenerate the electrons are
      (\S\ref{sec:ecthreshold}). What free-streaming does is \emph{prevent the
      drop from stopping}: if neutrinos were trapped they would build up a
      chemical potential until forward and reverse weak rates balanced, and
      \Ye{} would freeze at $\beta$ equilibrium. Escape removes that
      back-reaction, so the ratchet has nothing to push against. Hence \Ye{}
      declines essentially monotonically \textbf{through the burning phase in
      aggregate}, not term-by-term --- $\beta^-$ decays run the other way
      throughout, which is precisely why \S\ref{sec:softwired} treats
      \msn{80}'s \emph{absence} of $\beta$-decay partners as a structural
      defect rather than a harmless omission.
      \end{warn}
\end{enumerate}

That second point is worth pausing on. The three-row structure of $C$ ---
baryon, charge (with the electron column), lepton number ($e^-$ $+1$, $\nu$
$+1$, $\bar\nu$ $-1$) --- is not bookkeeping fussiness. It is what lets you
state ``lepton number is conserved by the reaction'' while the neutrino column
subsequently leaves the system.

\section{\texorpdfstring{$\langle E_\nu\rangle$}{<E nu>} --- an open item to
  understand before you meet it}
\label{sec:enu}

Energy loss per capture is not simply the $Q$-value: the emitted neutrino
carries a \emph{spectrum}, and what the zone loses is $\langle E_\nu\rangle$
per reaction. Weak-rate tables therefore provide both a rate and a neutrino
energy-loss rate --- \citeauthor{ffn}{} onward; \citet{langanke2001} tabulate effective rates
\emph{and} average neutrino energies, and \citet{patton2017} show exactly how
$\langle E_\nu\rangle$ is reconstructed from those tables --- and codes handle
the latter with varying care.

This is a \textbf{standing open item in the project} (``MESA
average-neutrino-energy handling ($\langle E_\nu\rangle$) for the neutrino
head''). When you design a neutrino output head, you will need to know exactly
what MESA reports and on what grid.

\begin{selfcheck}
\item Recompute $\lambda$ at $\rho = 10^{9}$ and $E_\nu = 5$~MeV. Is
      free-streaming still safe?
\item Explain, in terms of \S\ref{sec:freestreaming}, why carbon burning onward
      is neutrino-cooled while hydrogen burning is not.
\item The lepton-number row of $C$ assigns $\nu$ $+1$ and $\bar\nu$ $-1$. Write
      the row entries for an electron capture and for a $\beta^-$ decay, and
      verify each column closes.
\item Why can thermal neutrino losses be computed from $(T,\rho)$ alone, while
      weak neutrino losses cannot? Which one does the emulator have to learn?
\end{selfcheck}
